L'any 2013 es va celebrar el centenari del model atòmic proposat per Niels Bohr. Segons aquest model, l'àtom de $^{1}$H té un protó en el nucli i un electró que descriu una òrbita circular estable al seu voltant. El radi mínim que pot tenir aquesta òrbita, segons el model de Bohr, és de $5,29\times 10^{-11}\ \mathrm{m}$. Per a aquesta òrbita calculeu:

\figura[0.22]{fig1.png}
{\centering Niels Bohr\par}

\begin{apartat}{1}
La força elèctrica que actua sobre l'electró i la freqüència de gir que té.
\end{apartat}

\begin{apartat}{1}
L'energia mecànica de l'electró en l'òrbita que descriu al voltant del protó. Considereu negligible l'energia potencial gravitatòria.
\end{apartat}

\dades{%
  $k = 8,99\times 10^{9}\ \mathrm{N\,m^2\,C^{-2}}$\\
  $Q_\text{electró} = -1,60\times 10^{-19}\ \mathrm{C}$\\
  $m_\text{electró} = 9,11\times 10^{-31}\ \mathrm{kg}$\\
  $Q_\text{protó} = -Q_\text{electró}$\\
  $m_\text{protó} = 1,67\times 10^{-27}\ \mathrm{kg}$}
